\documentclass[]{article}

\usepackage{amsmath}
\usepackage{multirow}
\usepackage{natbib}
\usepackage{graphicx} 


\begin{document}

The transport of passive tracers by a turbulent flow is a fundamental process central to phenomena ranging from pollutant dispersal in geophysical flows to star formation in the interstellar medium. A primary goal of turbulence theory is to develop predictive models for this transport by connecting the statistical properties of particle trajectories to the underlying fluid dynamics. Such models often depend on assumptions about the flow's structure, particularly regarding the mechanisms of energy injection at large scales and the role of coherent structures. The geometry of the energy input can profoundly shape the flow's topology, creating distinct pathways for particle dispersion and potentially challenging classical assumptions of isotropic, homogeneous diffusion. Understanding this connection is therefore essential for accurately modeling transport in complex, multi-scale systems.

A key distinction in turbulence is the nature of the large-scale forcing, which can be broadly categorized as either compressive (irrotational) or solenoidal (rotational). Solenoidal forcing, which conserves volume and is characteristic of many subsonic and incompressible flows, naturally generates large-scale swirling motions and long-lived coherent vortices. These vortices can act as temporary transport barriers, trapping tracers and introducing correlations into their trajectories. This raises a central question: are these trapping events persistent enough to induce long-term memory and generate anomalous transport, where the Mean-Square Displacement grows faster than linearly with time? Furthermore, it remains unclear how the rotational nature of the forcing imposes a directional preference, or anisotropy, on particle dispersion relative to the large-scale flow field. Answering these questions is crucial for determining when classical diffusion models are applicable and how they must be modified to account for the specific physics of solenoidal turbulence.

In this work, we address these questions by analyzing the trajectories of thousands of passive Lagrangian tracers integrated within a high-resolution direct numerical simulation of subsonic, isothermal turbulence driven by purely solenoidal modes at the largest scales. This approach allows us to directly quantify both the persistence of vortex trapping and the anisotropy of the resulting dispersion. Our primary analysis focuses on the Mean-Square Displacement, which we decompose into components parallel and perpendicular to the instantaneous, local large-scale velocity vector. This decomposition provides a direct measure of directional transport. We complement this by identifying coherent vortices using the Q-criterion and explicitly measuring the residence times of tracers within these structures, allowing us to directly test the hypothesis that vortex trapping is the primary driver of anomalous transport.

Our results reveal two principal findings. First, we uncover a persistent and strong anisotropy in the dispersion process: transport perpendicular to the local large-scale velocity field systematically exceeds transport parallel to it. This transverse-dominant dispersion is a direct kinematic signature of the solenoidal forcing, where large-scale vortices sweep particles more efficiently in their rotational plane than along their axis. Second, we find that while tracers are transiently captured by coherent vortices, their residence times are remarkably short, lasting only a small fraction of a large-eddy turnover time. These brief trapping events, terminated by the rapid onset of three-dimensional vortex instabilities, are insufficient to generate the long-term memory required for sustained superdiffusion or Lévy flights. Consequently, the probability distributions of particle displacements remain nearly Gaussian at intermediate times before becoming light-tailed due to finite-domain effects. We conclude that in three-dimensional solenoidal turbulence, the forward energy cascade and inherent vortex dynamics effectively suppress long-range temporal correlations, ensuring an eventual return to a classical, albeit strongly anisotropic, diffusive regime.

\end{document}
                